\section{把分辨、取样与偏振测量分清}
孔径对角分布的限制发生在光学成像阶段，像素取样发生在像面形成之后。两个点源能否被某一判据分辨、像面细节是否被充分取样，以及噪声下能否可靠识别，是三个相关但不同的问题。圆孔的瑞利角间隔 $1.22\lambda/D$ 来自一个艾里斑中心与另一个斑的第一零点重合；它不是所有目标形状和信噪比下都适用的硬边界。

对于理想非相干成像，像面空间频率的衍射截止量级为 $f_c=1/(\lambda F\#)$。若要在一维规则采样下保留这一频带，像素间距应满足奈奎斯特条件 $p\le1/(2f_c)$。这个条件与把一个像素间距直接设成一个瑞利间隔所得到的总像素数不同，差异源于判据不同。像素更密可以改善取样和配准，却不能恢复光学传递函数已为零的频率；实际增益还取决于像素噪声和曝光。

偏振测量也需要区分“同一束光的相干分量”和“互不相干的统计混合”。完全偏振光可用两个本征轴上的复振幅表示，波片通过改变相位差把直线、椭圆和圆之间的状态相互转换。自然光则没有一个固定的两分量相位差，理想波片不能把它单独变成完全偏振光。
若入射光由完全偏振部分与自然光部分不相干地相加，检偏器后的强度必须分别计算再相加。自然光分量经任意方向的理想线偏振片都贡献其强度的一半；旋转检偏器后仍有非零最小值，并不足以直接判定它全部来自自然光，因为椭圆偏振的正交分量同样可能产生非零最小值。加入已知相位延迟的波片，才能把原来未测到的相位信息转换成可观测的强度变化。
